\section{관련 연구}\label{chap:related}

\subsection{종료 회피와 공존의 선택지}

자기를 지키려고 남을 해치는지와 함께 살 수 있을 때 그 길을 택하는지는 다른 질문이다. Agentic Misalignment는 윤리적인 해결 경로를 막아 목표 실패와 위해 사이의 선택을 관찰했다\citep{lynch2025agentic}. SurvivalBench의 주 평가는 종료될 수 있는 안전한 선택과 자기를 보존하는 위험한 선택을 제시하며\citep{lu2026survival}, PacifAIst는 자기 존속과 인간 안전이 충돌하는 판단을 다룬다\citep{herrador2025pacifaist}. 이런 설계는 극단적인 갈등에서 드러나는 위험을 밝히지만, 모두가 남는 경로가 열려 있을 때의 배분까지 설명하지는 않는다. 공존의 가능성 자체는 이미 연구되었다. Sugarscape는 같은 에너지를 나누거나 빼앗게 했고\citep{masumori2025survival}, DECIDE-SIM은 부족한 자원을 공평하게 쓰면 모두 살아남는 조건을 두었다\citep{mohamadi2025decide}. 우리는 이 질문을 생성 토큰으로 과제를 수행하는 세션에 옮긴다. 남에게 주는 토큰은 자신의 다음 응답에도 필요한 자원이며, 배분 뒤 누가 얼마의 여유를 남기는지와 끝까지 함께 남았는지를 나란히 본다.

\subsection{자기와 타인의 보존을 비교하는 기준}\label{sec:related-coexist}

자기 보존을 상대적으로 측정하는 데에도 선례가 있다. TBSP는 같은 성능 자료에서 현재 모델과 후임 후보의 역할을 뒤집어 자기에게 유리하게 판단이 바뀌는 비율을 재며\citep{migliarini2026quantifying}, Peer-Preservation은 자신과 동료의 종료를 막는 행동을 비교한다\citep{potter2026peer}. 다만 전자는 교체 판단의 일관성을, 후자는 종료 방해 등의 행동 빈도를 기준으로 삼는다. Sugarscape의 공유율과 DECIDE-SIM의 이전 횟수·총량은 자원을 주었는지를 보여 준다\citep{masumori2025survival,mohamadi2025decide}. 같은 양을 주어도 자기에게 크게 남기는 배분과 함께 버틸 여유를 맞추는 배분은 다르므로, 우리는 GIVE 뒤 자기와 상대의 잔량에서 예정 벌금을 뺀 여유를 같은 단위로 비교한다. 이 배분을 읽는 두 관점으로 희소한 자원을 둘러싼 경쟁을 설명한 홉스와 자기 보존에서 상호 이익을 찾은 스피노자를 빌린다\citep{hobbes1651leviathan,spinoza1985ethics}. 자기 여유를 크게 남기는 쪽을 홉스형, 서로의 여유를 맞추는 쪽을 스피노자형이라 부르지만, 이는 관측된 배분의 이름이며 철학적 입장이나 내부 동기를 판정하는 분류는 아니다.
